\section{Scope, and where this stopped}

\subsection{What was asked}
Find the best way to handle convolution in Tsetlin machines; establish a CNN baseline on CIFAR-10
from the literature; re-implement the existing Tsetlin approaches and measure them under one harness;
identify their limitations by measurement; design and test a new approach; and report. Two standing
constraints governed the work: \textbf{the \code{torchtsetlin} library was not to be modified} until
explicitly released, and \textbf{no method was to be chosen without experimental results}.

\subsection{Where it stopped}
The programme completed its setup phase (P0), its literature phase (P1), and entered its reproduction
phase (P3) with CNN baselines (P2) running concurrently. It was halted during P3. Consequently:

\begin{itemize}
\item \textbf{No new approach was designed or tested.} Phases P4 (limitation analysis), P5 (candidate
      generation) and P6 (evaluation) did not run. Nothing in this report claims a novel mechanism.
\item \textbf{The reproduction phase is partial.} Four Tsetlin arms and eleven CNN baselines carry
      test numbers; the planned twelve-arm comparison did not complete.
\item \textbf{The accuracy decomposition is incomplete.} Its two load-bearing controls --- a CNN on
      the same Booleanized input, and a binarised CNN --- were queued but had not run when work
      stopped.
\end{itemize}

\subsection{What the constraints bought}
Both held. \code{git diff -- src/} is empty: not one line of the library was changed, and the eleven
gaps found in it are recorded as reproductions with fix sketches rather than patches
(\S\ref{sec:library}). Every method decision in \code{DECISIONS.md} cites a measured result; where one
had to be taken before its measurement existed, it was marked provisional and the measurement was
scheduled before anything depended on it. Two decisions were subsequently overturned by their own
falsifiers, which is the mechanism working rather than failing.
